\section{Affine prime-successor loops}
\label{sec:affine-prime}

For an odd prime $p$ and $a\in\mathbb F_p\setminus\{0,1\}$, put
$b=1-a$. Define $Q_p(a)$ on $\mathbb F_p\cup\{\infty\}$ by
\[
 \infty*x=x*\infty=x,\qquad x*x=\infty\ (x\in\mathbb F_p),
 \qquad x*y=ax+by\ (x\ne y,\ x,y\in\mathbb F_p).
\]
In each finite row or column, the affine bijection fixes the line
label. Replacing that value by $\infty$ and placing the old label
in the new position proves that this is a Latin loop.
Write $[u]_p$ for the least positive residue of a nonzero field element.

\begin{theorem}[Three-view affine criterion]
\label{thm:affine-criterion}
The loop $Q_p(a)$ is FFF if and only if
\begin{equation}
\begin{split}
 &[a]_p,\ [a^{-1}]_p,\ [b^{-1}]_p\ \hbox{are even},\\
 &\operatorname{ord}_p(a),\operatorname{ord}_p(b),
 \operatorname{ord}_p(-b/a)\ \hbox{are even}.
\end{split}
 \label{eq:affine-criterion}
\end{equation}
\end{theorem}

\begin{proof}
Use a generic row parameter $t$, with $c=1-t$. A finite row
relative to the identity row consists of a two-cycle and
$(p-1)/\operatorname{ord}_p(c)$ cycles of length
$\operatorname{ord}_p(c)$. Translations and nonzero scalings of
finite labels are automorphisms and are two-transitive on them.
It remains to compare rows $0,1$. Their relative map
$P=r_1^{-1}r_0$, with $D=-t/c$, satisfies
\[
 P(\infty)=D,\quad P(0)=1,\quad P(c^{-1})=\infty,\quad
 P(y)=y+D\quad\hbox{at every other finite }y.
\]
Number the translation cycle by $jD$, $j\in\mathbb Z_p$, and
put $k=[-t^{-1}]_p$. Then $c^{-1}=kD$ and $1=(k+1)D$.
The redirected arrows give cycles of lengths $p-k$ and $k+1$,
both even exactly when $[t^{-1}]_p$ is even.

The column view has parameter $b$. A finite symbol $z$, viewed
as a map from columns to rows, swaps $z,\infty$ and otherwise
sends $y$ to $a^{-1}z-(b/a)y$. Its parameter is $a^{-1}$ and
its other coefficient is $-b/a$; the infinity symbol line is
the identity. Applying the row calculation to these three
parameters proves \eqref{eq:affine-criterion} in both directions.
\end{proof}

\begin{computationaltheorem}[A residue-class boundary control]
\label{cthm:fff98}
The explicitly defined loop $Q_{97}(36)$ is FFF of order $98$.
Thus $n\equiv2\pmod8$, $n>2$, is not a universal FFF obstruction.
\end{computationaltheorem}

Indeed, $b=62$, $36^{-1}=62$, $62^{-1}=36$, $-62/36=36$ in
$\mathbb F_{97}$, and $36^3=62^3=-1$. The orders and the three
least residues in \eqref{eq:affine-criterion} are even.
Primality and these identities are exact finite checks.
The materialized table also passes all
$3\binom{98}{2}=14259$ line-pair checks. This is not an
order-$18$ construction.

\subsection{A uniform large-prime bound}

\begin{theorem}[Eventual prime-successor existence]
\label{thm:affine-large-prime}
For every prime $p\geq10^{12}$, more than $p/1000$ parameters
produce FFF loops $Q_p(a)$.
\end{theorem}

The analytic input is Katz's published mixed-character-sum
estimate \cite[Theorem~1.5]{KatzMixed}. Its application and the
explicit cutoff are checked below, not inferred from finite tables.

\begin{proof}
Put $U=\mathbb F_p\setminus\{0,1\}$. For $p\equiv3\pmod4$,
take the quadratic character $\chi$ and the sufficient mask
\[
 W(a)=\tfrac14(1-\chi(a))(1-\chi(1-a)).
\]
For $p\equiv1\pmod4$, take a character of exact order four and use
\[
 W(a)=\tfrac18(1-\chi(a)^2)
       (1-\chi(1-a)+\chi(1-a)^2-\chi(1-a)^3).
\]
These are zero-one indicators. In the first case $a,b,-b/a$
are nonsquares. In the second, write $p-1=2^s m$ with $m$ odd.
The primitive-root exponent of $a$ is odd, that of $b$ is $2$
modulo $4$, and that of $-b/a$ is odd. All three multiplicative
orders are therefore even in either case. The mask's character
expansion has coefficient absolute sum one and constant coefficient
at least $1/8$.

Let $\psi(x)=\exp(2\pi i x/p)$ and let $\chi$ have order
$q\in\{2,4\}$. Except for the all-zero tuple, the sums
\[
 S(j,k;u,v,w)=
 \sum_{a\in U}\chi(a)^j\chi(1-a)^k
       \psi(ua+v/a+w/(1-a)),\quad 0\leq j,k<q,
\]
have absolute value at most $4\sqrt p$. For $u\ne0$, apply the
cited theorem with polynomial $ua$, rational remainder
$v/a+w/(1-a)$, and multiplicative polynomial $a^j(1-a)^k$.
All finite poles are simple; the pole order at infinity is below
one, and every nonpole zero has nontrivial character multiplicity.
With $t$ endpoint poles and $e$ other endpoint zeros, the estimate
is $(2t+e)\sqrt p$. Deleting the remaining $2-t-e$ endpoints
changes the sum by at most that number, so the result is at most
$4\sqrt p$. For $u=0,v\ne0$, substitute $z=1/a$; for
$u=v=0,w\ne0$, substitute $z=1/(1-a)$. Each is a bijection on
$U$, gives a nonzero linear phase plus at most one simple finite
pole, and changes the multiplicative factor to powers of $z,z-1$,
up to a constant. The same argument applies. With zero additive
phase, the nontrivial multiplicative sums are Jacobi sums or
one-character endpoint sums, bounded by $\sqrt p$.
The all-zero tuple has sum $p-2$.

Let $E(x)$ indicate nonzero even least residues. Geometric summation
gives its normalized Fourier coefficients
\[
 \widehat E(0)=\frac{p-1}{2p},\qquad
 |\widehat E(h)|=\frac{1}{2p|\cos(\pi h/p)|}\quad(h\ne0).
\]
Pairing $h,p-h$ and using $\sin x\geq2x/\pi$ on $[0,\pi/2]$
bounds the nonzero coefficient sum by
$\sum_{r=0}^{(p-3)/2}(2r+1)^{-1}\leq1+\tfrac12\log p$.
The full absolute sum is at most $L(p)=3/2+\tfrac12\log p$.
Expanding $W(a)E(a)E(a^{-1})E((1-a)^{-1})$ therefore bounds
its parameter count by
\begin{equation}
 N_p\geq\frac{(p-2)(p-1)^3}{64p^3}
                   -4\sqrt p\,L(p)^3.
 \label{eq:large-prime-count}
\end{equation}
The first term is at least $(p-5)/64\geq p/65$ for $p\geq325$.
The function $\sqrt x/L(x)^3$ increases once $L(x)>3$.
At $x=10^{12}$, $\log10<2.303$ gives $L(x)<15.318$ and
$260(15.318)^3<935000<10^6$, exact rational inequalities.
It follows that $N_p>(1-0.935)p/65=p/1000$.
For example, the finite rational inequality
$\sum_{k=0}^9(2303/1000)^k/k!>10$ proves the logarithm bound.
\end{proof}

For $p>3$ these loops are not group-isotopic. Associativity at
$(0,0,1)$ would force $b^2=1$, hence $b=-1$, and at $(1,0,0)$
would force $a=-1$, contradicting $a+b=1$. A group-isotopic loop
is a group by the basis-family closure criterion: since its identity
row is the identity permutation, closure gives $L_xL_y=L_{xy}$
after evaluating at the identity, hence associativity.

\subsection{Exact classification inside the affine family}

\begin{theorem}[Affine-family main classes]
\label{thm:affine-anharmonic}
For $p\geq5$, the loops $Q_p(a),Q_p(a')$ are main-class-equivalent
exactly when
\[
 a'\in\mathcal A(a)=
 \{a,1-a,1/a,1/(1-a),a/(a-1),(a-1)/a\}.
\]
Distinct parameters are never isomorphic; each loop has automorphism
group $\operatorname{AGL}(1,p)$.
\end{theorem}

\begin{proof}
Swapping the inputs and solving the relation for either input give
the displayed parameters, including all triples involving
$\infty$. Thus these relations suffice.

The types of row pairs involving infinity and of two finite rows are
\[
 (2,\operatorname{ord}_p(b)^{(p-1)/\operatorname{ord}_p(b)}),
 \qquad([a^{-1}]_p,p+1-[a^{-1}]_p).
\]
They coincide only if $\operatorname{ord}_p(b)=p-1$ and
$[a^{-1}]_p\in\{2,p-1\}$. The column version exchanges $a,b$.
Outside $E=\{2,1/2,-1\}$ both identity lines are intrinsic:
each is the unique line with all $p$ incident pair types the
same, whereas a finite line has both types. An isotopy between
such loops must fix the two identity lines; the loop equations
make its three maps equal. Its finite-field bijection $f$ satisfies
\[
 f(ax+by)=a'f(x)+(1-a')f(y),
\]
also at $x=y$. Subtract $f(0)$ and set one argument to zero.
Inverting $a,b$ in the resulting identities proves additivity of
$h(x)=f(x)-f(0)$. On a prime field $h$ is multiplication by a
nonzero scalar; the identity then forces $a=a'$.

For a paratopy first perform its coordinate permutation, replacing
$a$ by $t\in\mathcal A(a)$. The set $E$ is a single orbit.
If exactly one of $t,a'$ is in $E$, either its special and finite
pair types coincide in some view, giving an invariant mismatch,
or both identity lines remain intrinsic and the same argument
forces equality. If both are in $E$, they already lie in the
same displayed orbit. This proves necessity.
For an isomorphism the identity is fixed without exceptions,
so the additive argument always gives $a=a'$ and $f(x)=ux+v$.
Conversely every such map with $u\ne0$, fixing $\infty$, is
an automorphism.
\end{proof}

Let $S_p$ be the full parameter set satisfying
\eqref{eq:affine-criterion}. The anharmonic orbits have size six,
apart from $E$ of size three and the two roots of $a^2-a+1=0$
when present. Let $e_p,h_p$ indicate whether these exceptional
orbits lie in $S_p$. The exact number of represented FFF main
classes is
\[
 C_p^{\mathrm{aff}}=(|S_p|+3e_p+4h_p)/6.
\]
Indeed, the three involutions fix only $1/2,2,-1$, and the
order-three elements fix exactly the two polynomial roots.
For $p\geq5$ these exceptional sets are disjoint. Consequently,
Theorem~\ref{thm:affine-large-prime} gives more than $p/6000$
FFF main classes at every prime-successor order $p+1$ with
$p\geq10^{12}$.

Finite controls construct all $253$ parameters for the $13$ odd
primes through $43$, compare all views with the criterion, and
check parameter transformations and orbit counts. They do not
prove the analytic estimate or its uniform cutoff. At $p=17$
the criterion excludes only the scalar affine family, not all
order-$18$ squares. The classification is within this family,
not a census of all Latin squares. Literature priority and external
expert review of the analytic application remain open.
